% Options for packages loaded elsewhere
\PassOptionsToPackage{unicode}{hyperref}
\PassOptionsToPackage{hyphens}{url}
\PassOptionsToPackage{space}{xeCJK}
\documentclass[
  a4paper,11pt]{ltjsarticle}
\usepackage{xcolor}
\usepackage[margin=1in]{geometry}
\usepackage{amsmath,amssymb}
\usepackage{graphicx}
\setcounter{secnumdepth}{-\maxdimen} % remove section numbering
\usepackage{iftex}
\ifPDFTeX
  \usepackage[T1]{fontenc}
  \usepackage[utf8]{inputenc}
  \usepackage{textcomp}
\else
  \usepackage{unicode-math}
  \defaultfontfeatures{Scale=MatchLowercase}
  \defaultfontfeatures[\rmfamily]{Ligatures=TeX,Scale=1}
\fi
\usepackage{lmodern}
\ifPDFTeX\else
  \setmainfont[]{Times New Roman}
  \ifXeTeX
    \usepackage{xeCJK}
    \setCJKmainfont[]{Hiragino Mincho ProN}
  \fi
  \ifLuaTeX
    \usepackage[]{luatexja-fontspec}
    \setmainjfont[]{Hiragino Mincho ProN}
  \fi
\fi
\IfFileExists{upquote.sty}{\usepackage{upquote}}{}
\IfFileExists{microtype.sty}{%
  \usepackage[]{microtype}
  \UseMicrotypeSet[protrusion]{basicmath}
}{}
\makeatletter
\@ifundefined{KOMAClassName}{%
  \IfFileExists{parskip.sty}{%
    \usepackage{parskip}
  }{%
    \setlength{\parindent}{0pt}
    \setlength{\parskip}{6pt plus 2pt minus 1pt}}
}{% if KOMA class
  \KOMAoptions{parskip=half}}
\makeatother
\usepackage{longtable,booktabs,array}
\newcounter{none}
\usepackage{calc}
\usepackage{etoolbox}
\makeatletter
\patchcmd\longtable{\par}{\if@noskipsec\mbox{}\fi\par}{}{}
\makeatother
\IfFileExists{footnotehyper.sty}{\usepackage{footnotehyper}}{\usepackage{footnote}}
\makesavenoteenv{longtable}
\usepackage{bookmark}
\IfFileExists{xurl.sty}{\usepackage{xurl}}{}
\urlstyle{same}
\usepackage{hyperref}
\hypersetup{
  hidelinks,
  pdfcreator={LaTeX via pandoc}}
\setlength{\emergencystretch}{3em}
\providecommand{\tightlist}{%
  \setlength{\itemsep}{0pt}\setlength{\parskip}{0pt}}
\usepackage{amsthm}
\newtheorem{theorem}{定理}[section]
\newtheorem{proposition}[theorem]{命題}
\newtheorem{lemma}[theorem]{補題}
\newtheorem{corollary}[theorem]{系}
\theoremstyle{definition}
\newtheorem{definition}[theorem]{定義}
\theoremstyle{remark}
\newtheorem{remark}[theorem]{注意}
\newtheorem{observation}[theorem]{観察}
\begin{document}
\section{6次元超直方体と中心投影との関係性の考察}\label{ux6b21ux5143ux8d85ux76f4ux65b9ux4f53ux3068ux4e2dux5fc3ux6295ux5f71ux3068ux306eux95a2ux4fc2ux6027ux306eux8003ux5bdf}

\subsection{------ 3つの曲率半径 R, R₀, R₁ の構造
------}\label{ux3064ux306eux66f2ux7387ux534aux5f84-r-rux2080-rux2081-ux306eux69cbux9020}

\textbf{著者}: 木原 範昭 (WF System Co., Ltd.) \textbf{日付}: 2026年4月

\begin{center}\rule{0.5\linewidth}{0.5pt}\end{center}

\subsection{0. 本稿の目的}\label{ux672cux7a3fux306eux76eeux7684}

\begin{center}\rule{0.5\linewidth}{0.5pt}\end{center}

\subsection{1.
準備：本論文での前提条件と仮定}\label{ux6e96ux5099ux672cux8ad6ux6587ux3067ux306eux524dux63d0ux6761ux4ef6ux3068ux4eeeux5b9a}

\subsubsection{1.0
単位系の宣言}\label{ux5358ux4f4dux7cfbux306eux5ba3ux8a00}

本稿では、宇宙論的スケール（曲率半径
\(R_0\)）と素粒子スケール（陽子電荷半径
\(R_1\)）を同一の単位系で統一的に扱うため、\textbf{幾何学的単位系}（\(c = 1\),
\(G = 1\)）を採用する。

\[c = 1, \quad G = 1\]

この単位系では \([\text{長さ}] = [\text{時間}] = [\text{質量}]\)
となり、すべてのスケールを長さ（メートル）で統一的に比較できる。

\textbf{\(\hbar = 1\) は採用しない。}
粒子物理で慣用される自然単位系（\(c = \hbar = 1\)）では長さと
\(1/\text{エネルギー}\)
が同一視され、宇宙論的スケールと素粒子スケールの比 \(R_0 / R_1\)
の幾何学的意味が不透明になる。本稿の目的は両スケールの比を幾何学的に定式化することにあるため、\(\hbar\)
は物理定数として陽に保持する。

\subsubsection{\texorpdfstring{1.1 空間曲率 \(R_0\)
の仮定}{1.1 空間曲率 R\_0 の仮定}}\label{ux7a7aux9593ux66f2ux7387-r_0-ux306eux4eeeux5b9a}

現代宇宙論において、宇宙の曲率を特徴づける量には複数の定義がある。以下に、現在の観測値（Planck
2018 {[}5{]}）に基づく曲率半径を、\(c = 1\) の幾何学的単位系で示す。

{\def\LTcaptype{none} % do not increment counter
\begin{longtable}[]{@{}
  >{\raggedright\arraybackslash}p{(\linewidth - 4\tabcolsep) * \real{0.3333}}
  >{\raggedright\arraybackslash}p{(\linewidth - 4\tabcolsep) * \real{0.3333}}
  >{\raggedright\arraybackslash}p{(\linewidth - 4\tabcolsep) * \real{0.3333}}@{}}
\toprule\noalign{}
\begin{minipage}[b]{\linewidth}\raggedright
定義
\end{minipage} & \begin{minipage}[b]{\linewidth}\raggedright
曲率半径
\end{minipage} & \begin{minipage}[b]{\linewidth}\raggedright
値（\(c=1\)）
\end{minipage} \\
\midrule\noalign{}
\endhead
\bottomrule\noalign{}
\endlastfoot
\(\Lambda\) 由来 & \(R_\Lambda = \sqrt{3/\Lambda}\) &
\(5.5 \times 10^{17}\)（174億光年） \\
ハッブル半径 & \(R_H = c/H_0\) & \(4.6 \times 10^{17}\)（145億光年） \\
空間曲率 & \(R_K = 1/\sqrt{\lvert K \rvert}\) &
\(> 7.3 \times 10^{18}\)（\textgreater{} 2300億光年） \\
\end{longtable}
}

本稿では、空間曲率 \(R_K\) を \(R_0\)
として採用する。以下にその算出根拠を述べる。

\paragraph{\texorpdfstring{空間曲率 \(R_K\)
の算出根拠}{空間曲率 R\_K の算出根拠}}\label{ux7a7aux9593ux66f2ux7387-r_k-ux306eux7b97ux51faux6839ux62e0}

FLRW計量における空間のガウス曲率 \(K\) は、密度パラメータ \(\Omega_k\)
とハッブル定数 \(H_0\) により

\[K = -\,\Omega_k \frac{H_0^2}{c^2}\]

で定義される。ここで \(\Omega_k = 1 - \Omega_{\mathrm{total}}\)
は全エネルギー密度の臨界密度からのずれを表す無次元量であり、Planck 2018
の観測から

\[\Omega_k = 0.001 \pm 0.002 \quad (68\% \text{ CL})\]

と報告されている。中心値は
\(\Omega_k = 0\)（空間的に平坦）と矛盾せず、現在の観測精度では曲率は検出されていない。

\(2\sigma\) の上限 \(\lvert \Omega_k \rvert < 0.004\)
を用いると、ガウス曲率の上限は

\[\lvert K \rvert < 0.004 \times \frac{H_0^2}{c^2} = 0.004 \times 5.31 \times 10^{-53} \text{ m}^{-2} = 2.1 \times 10^{-55} \text{ m}^{-2}\]

となる。対応する曲率半径の下限は

\[R_K = \frac{1}{\sqrt{\lvert K \rvert}} > \frac{1}{\sqrt{2.1 \times 10^{-55}}} = 2.2 \times 10^{27} \text{ m}\]

\(c = 1\) に換算すると

\[R_0 \equiv R_K > \frac{2.2 \times 10^{27}}{c} = \frac{2.2 \times 10^{27}}{3.0 \times 10^{8}} = 7.3 \times 10^{18}\]

すなわち
2300億光年に相当する。これは観測可能な宇宙の半径（約465億光年）の約50倍であり、空間は局所的にはほぼ平坦であるが、有限の曲率半径を持つ可能性は排除されていない。

本稿では、この下限値を仮定として

\[R_0 = 7.3 \times 10^{18} \quad (c = 1)\]

と置く。

\subsubsection{\texorpdfstring{1.2 素粒子曲率 \(R_1\)
の仮定}{1.2 素粒子曲率 R\_1 の仮定}}\label{ux7d20ux7c92ux5b50ux66f2ux7387-r_1-ux306eux4eeeux5b9a}

現代素粒子論において、素粒子の「大きさ」を特徴づける量にも複数の定義がある。以下に、代表的な粒子について現在の実験値を示す。

{\def\LTcaptype{none} % do not increment counter
\begin{longtable}[]{@{}
  >{\raggedright\arraybackslash}p{(\linewidth - 6\tabcolsep) * \real{0.2500}}
  >{\raggedright\arraybackslash}p{(\linewidth - 6\tabcolsep) * \real{0.2500}}
  >{\raggedright\arraybackslash}p{(\linewidth - 6\tabcolsep) * \real{0.2500}}
  >{\raggedright\arraybackslash}p{(\linewidth - 6\tabcolsep) * \real{0.2500}}@{}}
\toprule\noalign{}
\begin{minipage}[b]{\linewidth}\raggedright
粒子
\end{minipage} & \begin{minipage}[b]{\linewidth}\raggedright
半径の定義
\end{minipage} & \begin{minipage}[b]{\linewidth}\raggedright
SI値
\end{minipage} & \begin{minipage}[b]{\linewidth}\raggedright
\(c=1\)
\end{minipage} \\
\midrule\noalign{}
\endhead
\bottomrule\noalign{}
\endlastfoot
陽子 & 電荷半径 & \(8.41 \times 10^{-16}\) m（0.841 fm） &
\(2.81 \times 10^{-24}\) \\
中性子 & 電荷半径 &
\(\sqrt{\langle r^2 \rangle} \approx 8.4 \times 10^{-16}\) m &
\(2.80 \times 10^{-24}\) \\
荷電パイ中間子 & 電荷半径 & \(6.6 \times 10^{-16}\) m（0.66 fm） &
\(2.20 \times 10^{-24}\) \\
電子 & 古典電子半径 & \(2.82 \times 10^{-15}\) m &
\(9.40 \times 10^{-24}\) \\
電子 & コンプトン波長 \(\bar{\lambda}_e = \hbar / m_e c\) &
\(3.86 \times 10^{-13}\) m & \(1.29 \times 10^{-21}\) \\
電子 & 実験上限（点粒子） & \(< 10^{-18}\) m &
\(< 3.3 \times 10^{-27}\) \\
光子 & 実験上限（点粒子） & \(< 10^{-18}\) m &
\(< 3.3 \times 10^{-27}\) \\
ニュートリノ & 実験上限（点粒子） & \(< 10^{-18}\) m &
\(< 3.3 \times 10^{-27}\) \\
\end{longtable}
}

\textbf{注記:} - 陽子電荷半径は PRad (2019) {[}6{]} および CODATA 2018
の値 \(r_p = 0.8414 \pm 0.0019\) fm に基づく。 -
電子・光子・ニュートリノは現在の実験精度では内部構造が検出されておらず、点粒子として扱われる。
- 古典電子半径 \(r_e = e^2 / (4\pi\varepsilon_0 m_e c^2)\)
は散乱断面積のスケールであり、粒子の「大きさ」とは異なる。

本稿では、実験的に内部構造が確認されている最も基本的な粒子として
\textbf{陽子の電荷半径} を \(R_1\) として採用する。

\[R_1 = 2.81 \times 10^{-24} \quad (c = 1)\]

\begin{center}\rule{0.5\linewidth}{0.5pt}\end{center}

\subsection{2.
素粒子のモデル化}\label{ux7d20ux7c92ux5b50ux306eux30e2ux30c7ux30ebux5316}

\textbf{集合符号の幾何学的解釈について.} 論文 {[}4{]} では集合符号 \(0\)
を「座標値
\(= 0\)」と定義したが、これは粒子がその軸上の原点に位置するという意味ではない。幾何学的には、集合符号
\(0\)
は\textbf{その軸方向の成分を持たない}こと、すなわち\textbf{その軸と直交している}ことを意味する。したがって、ある粒子の集合符号がある軸について
\(0\)
であるとは、その粒子の状態がその軸の値に依存しないということである。

\textbf{集合符号の物理的解釈の限界について.}
本章で列挙する各粒子の集合符号は、論文 {[}4{]}
で定義された組合せ論的な構造であり、各軸の非零成分（\(+\), \(-\),
\(\pm\)）が物理的に何を意味するか------すなわち、その成分がどのような物理量に対応し、どのような値を取るか------は、現時点では未解決の問題である。以下では集合符号の構造的特徴のみを記述し、物理的解釈は後の章で検討する。

\subsubsection{2.1
光子の6次元超直方体のモデル化}\label{ux5149ux5b50ux306e6ux6b21ux5143ux8d85ux76f4ux65b9ux4f53ux306eux30e2ux30c7ux30ebux5316}

論文 {[}4{]} §10.10 において、光子 \(\gamma\)
は以下の集合符号で定義される。

\[\gamma: \quad (x, y, z, t, R, Q) = (0, 0, 0, 0, +, 0)\]

各軸の意味を明示すると：

{\def\LTcaptype{none} % do not increment counter
\begin{longtable}[]{@{}ccl@{}}
\toprule\noalign{}
軸 & 集合符号 & 意味 \\
\midrule\noalign{}
\endhead
\bottomrule\noalign{}
\endlastfoot
\(x\) & \(0\) & 空間第1軸と直交：\(x\) 軸の値に依存しない \\
\(y\) & \(0\) & 空間第2軸と直交：\(y\) 軸の値に依存しない \\
\(z\) & \(0\) & 空間第3軸と直交：\(z\) 軸の値に依存しない \\
\(t\) & \(0\) & 時間軸と直交：\(t\) 軸の値に依存しない \\
\(R\) & \(+\) & 曲率半径軸：正の成分を持つ（\(R > 0\)） \\
\(Q\) & \(0\) & 色荷ラベル：\(Q = 0\)（色中性） \\
\end{longtable}
}

光子はスピン1ボソン（\(n = 1\)、\(R\) 軸のみ非零）であり、\(x, y, z, t\)
軸のいずれとも直交している（これらの軸の値に依存しない）。\(Q = 0\)（色中性）である。

ここで注目すべきは、光子を特徴づける唯一の非零成分が \(R\)
軸であり、その集合符号が「\(+\)」、すなわち \(R > 0\)
であることである。論文 {[}4{]} の集合符号の定義（§1.3.3）では、\(+\)
は座標値が正であることのみを要求し、\textbf{具体的な値は指定されない}。

すなわち、\(R\) 軸の値は \(R > 0\)
を満たす任意の正の実数であり、光子の6次元超直方体モデルにおいて \(R\)
は\textbf{唯一の自由パラメータ}である。

§1.1 で定めた宇宙論的曲率半径 \(R_0 = 7.3 \times 10^{18}\) と、§1.2
で定めた陽子電荷半径 \(R_1 = 2.81 \times 10^{-24}\) はいずれも \(R > 0\)
の範囲にあるが、その値は42桁異なる。光子の \(R\)
にどの値を仮定するかは、本モデルにおける未解決の問題である。

\subsubsection{2.2
グラビトンの6次元超直方体のモデル化}\label{ux30b0ux30e9ux30d3ux30c8ux30f3ux306e6ux6b21ux5143ux8d85ux76f4ux65b9ux4f53ux306eux30e2ux30c7ux30ebux5316}

論文 {[}4{]} §10.10 において、グラビトン \(G\)
は以下の集合符号で定義される。

\[G: \quad (x, y, z, t, R, Q) = (0, 0, 0, \pm, \pm, 0)\]

各軸の意味を明示すると：

{\def\LTcaptype{none} % do not increment counter
\begin{longtable}[]{@{}ccl@{}}
\toprule\noalign{}
軸 & 集合符号 & 意味 \\
\midrule\noalign{}
\endhead
\bottomrule\noalign{}
\endlastfoot
\(x\) & \(0\) & 空間第1軸と直交：\(x\) 軸の値に依存しない \\
\(y\) & \(0\) & 空間第2軸と直交：\(y\) 軸の値に依存しない \\
\(z\) & \(0\) & 空間第3軸と直交：\(z\) 軸の値に依存しない \\
\(t\) & \(\pm\) & 時間軸：非零の成分を持つ（\(t \neq 0\)） \\
\(R\) & \(\pm\) & 曲率半径軸：非零の成分を持つ（\(R \neq 0\)） \\
\(Q\) & \(0\) & 色荷ラベル：\(Q = 0\)（色中性） \\
\end{longtable}
}

グラビトンはスピン2ボソン（\(n = 2\)、\(t\) 軸と \(R\)
軸の2本が非零）であり、\(x, y, z\)
軸のいずれとも直交している。\(Q = 0\)（色中性）である。

光子との比較において注目すべき点は、グラビトンが \(t\) 軸と \(R\)
軸の\textbf{両方}に非零成分を持つことである。光子が \(R\)
軸のみを非零成分として持つのに対し、グラビトンは \(t\) 軸と \(R\)
軸が結合した構造を持つ。すなわち、重力の媒介粒子は時間と曲率半径の双方に関与する。

また、光子と同様に \(R\)
軸の集合符号は「\(\pm\)」（\(R \neq 0\)）であり、\(R\)
の具体的な値は指定されない。

\subsubsection{2.3
ヒッグス粒子の6次元超直方体のモデル化}\label{ux30d2ux30c3ux30b0ux30b9ux7c92ux5b50ux306e6ux6b21ux5143ux8d85ux76f4ux65b9ux4f53ux306eux30e2ux30c7ux30ebux5316}

論文 {[}4{]} §10.10 において、ヒッグス粒子 \(H^0\)
は以下の集合符号で定義される。

\[H^0: \quad (x, y, z, t, R, Q) = (\pm, 0, 0, 0, 0, 0)\]

各軸の意味を明示すると：

{\def\LTcaptype{none} % do not increment counter
\begin{longtable}[]{@{}ccl@{}}
\toprule\noalign{}
軸 & 集合符号 & 意味 \\
\midrule\noalign{}
\endhead
\bottomrule\noalign{}
\endlastfoot
\(x\) & \(\pm\) & 空間第1軸：非零の成分を持つ（\(x \neq 0\)） \\
\(y\) & \(0\) & 空間第2軸と直交：\(y\) 軸の値に依存しない \\
\(z\) & \(0\) & 空間第3軸と直交：\(z\) 軸の値に依存しない \\
\(t\) & \(0\) & 時間軸と直交：\(t\) 軸の値に依存しない \\
\(R\) & \(0\) & 曲率半径軸と直交：\(R\) 軸の値に依存しない \\
\(Q\) & \(0\) & 色荷ラベル：\(Q = 0\)（色中性） \\
\end{longtable}
}

ヒッグス粒子はスピン0ボソン（\(n = 0\)、\(t, R, Q\)
軸すべてと直交）であり、\(x, y, z\)
軸のうち1軸にのみ符号付き長さ成分（\(\pm\)）を持つ。\(x, y, z\)
軸は対称であるから（論文 {[}4{]} 命題
1.1）、どの1軸が非零であるかは表記ラベルの選択に過ぎない。

\subsubsection{2.4
電子の6次元超直方体のモデル化}\label{ux96fbux5b50ux306e6ux6b21ux5143ux8d85ux76f4ux65b9ux4f53ux306eux30e2ux30c7ux30ebux5316}

論文 {[}4{]} §10.10 において、電子 \(e^-\)
は以下の集合符号で定義される。

\[e^-: \quad (x, y, z, t, R, Q) = (+, +, 0, 0, 0, 0)\]

各軸の意味を明示すると：

{\def\LTcaptype{none} % do not increment counter
\begin{longtable}[]{@{}ccl@{}}
\toprule\noalign{}
軸 & 集合符号 & 意味 \\
\midrule\noalign{}
\endhead
\bottomrule\noalign{}
\endlastfoot
\(x\) & \(+\) & 空間第1軸：正の成分を持つ（\(x > 0\)） \\
\(y\) & \(+\) & 空間第2軸：正の成分を持つ（\(y > 0\)） \\
\(z\) & \(0\) & 空間第3軸と直交：\(z\) 軸の値に依存しない \\
\(t\) & \(0\) & 時間軸と直交：\(t\) 軸の値に依存しない \\
\(R\) & \(0\) & 曲率半径軸と直交：\(R\) 軸の値に依存しない \\
\(Q\) & \(0\) & 色荷ラベル：\(Q = 0\)（色中性） \\
\end{longtable}
}

電子はスピン1/2フェルミオン（\(xyz\)
軸のうち2本が非零、\(n = 0\)）であり、\(t, R\)
軸と直交し、\(Q = 0\)（レプトン、色中性）、第1世代（\(x, y\)
の2軸が非零）である。\(x = +\), \(y = +\) であるから、電子は \(xy\)
平面上に符号なし面積（正×正 = 正）を持つ。

\subsubsection{2.5
陽子の6次元超直方体のモデル化}\label{ux967dux5b50ux306e6ux6b21ux5143ux8d85ux76f4ux65b9ux4f53ux306eux30e2ux30c7ux30ebux5316}

陽子 \(p\) は複合粒子であり、3つのクォーク \(u_r, u_g, d_b\)
から構成される。論文 {[}4{]} §10.10
に基づき、各構成クォークの集合符号を列挙する。

\textbf{アップクォーク \(u_r\)（色荷 \(r\)）：}

\[u_r: \quad (x, y, z, t, R, Q) = (+, +, 0, 0, 0, r)\]

{\def\LTcaptype{none} % do not increment counter
\begin{longtable}[]{@{}ccl@{}}
\toprule\noalign{}
軸 & 集合符号 & 意味 \\
\midrule\noalign{}
\endhead
\bottomrule\noalign{}
\endlastfoot
\(x\) & \(+\) & 空間第1軸：正の成分を持つ（\(x > 0\)） \\
\(y\) & \(+\) & 空間第2軸：正の成分を持つ（\(y > 0\)） \\
\(z\) & \(0\) & 空間第3軸と直交：\(z\) 軸の値に依存しない \\
\(t\) & \(0\) & 時間軸と直交：\(t\) 軸の値に依存しない \\
\(R\) & \(0\) & 曲率半径軸と直交：\(R\) 軸の値に依存しない \\
\(Q\) & \(r\) & 色荷ラベル：赤 \\
\end{longtable}
}

\textbf{アップクォーク \(u_g\)（色荷 \(g\)）：}

\[u_g: \quad (x, y, z, t, R, Q) = (+, +, 0, 0, 0, g)\]

{\def\LTcaptype{none} % do not increment counter
\begin{longtable}[]{@{}ccl@{}}
\toprule\noalign{}
軸 & 集合符号 & 意味 \\
\midrule\noalign{}
\endhead
\bottomrule\noalign{}
\endlastfoot
\(x\) & \(+\) & 空間第1軸：正の成分を持つ（\(x > 0\)） \\
\(y\) & \(+\) & 空間第2軸：正の成分を持つ（\(y > 0\)） \\
\(z\) & \(0\) & 空間第3軸と直交：\(z\) 軸の値に依存しない \\
\(t\) & \(0\) & 時間軸と直交：\(t\) 軸の値に依存しない \\
\(R\) & \(0\) & 曲率半径軸と直交：\(R\) 軸の値に依存しない \\
\(Q\) & \(g\) & 色荷ラベル：緑 \\
\end{longtable}
}

\textbf{ダウンクォーク \(d_b\)（色荷 \(b\)）：}

\[d_b: \quad (x, y, z, t, R, Q) = (+, -, 0, 0, 0, b)\]

{\def\LTcaptype{none} % do not increment counter
\begin{longtable}[]{@{}ccl@{}}
\toprule\noalign{}
軸 & 集合符号 & 意味 \\
\midrule\noalign{}
\endhead
\bottomrule\noalign{}
\endlastfoot
\(x\) & \(+\) & 空間第1軸：正の成分を持つ（\(x > 0\)） \\
\(y\) & \(-\) & 空間第2軸：負の成分を持つ（\(y < 0\)） \\
\(z\) & \(0\) & 空間第3軸と直交：\(z\) 軸の値に依存しない \\
\(t\) & \(0\) & 時間軸と直交：\(t\) 軸の値に依存しない \\
\(R\) & \(0\) & 曲率半径軸と直交：\(R\) 軸の値に依存しない \\
\(Q\) & \(b\) & 色荷ラベル：青 \\
\end{longtable}
}

3つのクォークはいずれも \(R\) 軸と直交している（\(R\)
軸の値に依存しない）。

面積の符号に着目すると、アップクォーク \(u\) は \(x = +\), \(y = +\)
であるから \(xy\) 平面上に符号なし面積（正×正 =
正）を持ち、ダウンクォーク \(d\) は \(x = +\), \(y = -\)
であるから符号付き面積（正×負 = 負）を持つ。陽子 \(p = u_r + u_g + d_b\)
の3つの構成クォークのうち、ダウンクォークのみが負の面積を持つ。

\subsubsection{\texorpdfstring{2.6 \(R\)
軸に着目した比較}{2.6 R 軸に着目した比較}}\label{r-ux8ef8ux306bux7740ux76eeux3057ux305fux6bd4ux8f03}

ここまでの粒子を \(R\) 軸に着目して比較すると：

{\def\LTcaptype{none} % do not increment counter
\begin{longtable}[]{@{}lcccl@{}}
\toprule\noalign{}
粒子 & スピン & \(R\) 軸 & \(t\) 軸 & 空間軸 \\
\midrule\noalign{}
\endhead
\bottomrule\noalign{}
\endlastfoot
光子 \(\gamma\) & 1 & \(+\) & \(0\) & すべて \(0\) \\
グラビトン \(G\) & 2 & \(\pm\) & \(\pm\) & すべて \(0\) \\
ヒッグス \(H^0\) & 0 & \(0\) & \(0\) & \(x = \pm\) \\
電子 \(e^-\) & 1/2 & \(0\) & \(0\) & \(x = +, y = +\) \\
クォーク \(u, d\) & 1/2 & \(0\) & \(0\) & \(x = +, y = \pm\) \\
\end{longtable}
}

\(R\)
軸に成分を持つ粒子（光子、グラビトン）はいずれも質量ゼロであり、空間軸すべてと直交している。一方、\(R\)
軸と直交する粒子（ヒッグス、電子、クォーク）は質量を持ち、空間軸に非零成分を持つ。

\begin{center}\rule{0.5\linewidth}{0.5pt}\end{center}

\subsection{3.}\label{section}

\begin{center}\rule{0.5\linewidth}{0.5pt}\end{center}

\subsection{参考文献}\label{ux53c2ux8003ux6587ux732e}

{[}1{]} 木原範昭, 「グノモン正写像による4次元時空の幾何学的定式化」,
2025. {[}2{]} 木原範昭,
「中心投影による主観空間の構成と位相空間における相互作用の定式化」,
2026. {[}3{]} 木原範昭,
「離散空間における中心投影の全球被覆と5次元背景空間の整数論的必然性」,
2026. {[}4{]} 木原範昭, 「6次元超直方体の集合構造とその組合せ論的性質」,
2026.
\end{document}
